% Figure for Oscillations, Waves and Optics §A1.2: the Lennard-Jones potential and its harmonic approximation at the minimum (Example §A1.2.2), in units of sigma and epsilon.
\begin{tikzpicture}[font=\small]
  \begin{axis}[nfig, width=9cm, height=5.4cm, xmin=0.9, xmax=2.6, ymin=-1.25, ymax=1.0,
      axis x line=middle, axis y line=left, clip=true,
      xtick={1,1.1225,2}, xticklabels={$\sigma$,,$2\sigma$}, ytick={-1}, yticklabels={$-\varepsilon$},
      extra x ticks={1.1225}, extra x tick labels={$x_{\min}$}, extra x tick style={tick label style={above, yshift=2pt}},
      xlabel={$x$}, xlabel style={anchor=west}, ylabel={$U$}, ylabel style={rotate=-90}]
    \addplot[figB, domain=0.96:2.6, samples=400] {4*((1/x)^12-(1/x)^6)};
    \addplot[figR, dashed, domain=0.9:1.5, samples=100] {-1+0.5*57.146*(x-1.1225)^2};
    \draw[black!40, thin, densely dotted] (axis cs:1.1225,0) -- (axis cs:1.1225,-1);
    \draw[black!40, thin, densely dotted] (axis cs:0.9,-1) -- (axis cs:2.6,-1);
  \end{axis}
\end{tikzpicture}
